\documentclass[journal]{IEEEtran}
\usepackage{amsmath,amsfonts,amssymb}
\usepackage{cite}
\usepackage{graphicx}
\usepackage{booktabs}
\usepackage{multirow}
\usepackage{array}
\usepackage{url}
\usepackage{xcolor}

\begin{document}

\title{LITHYX: Physics-Informed, Calibrated Uncertainty Digital Twin Framework for Battery Prognostics and Lifecycle Decision Support}

\author{Balasurya~et~al.,~\IEEEmembership{Member,~IEEE}
\thanks{Manuscript received October 3, 2026. Corresponding author: Balasurya.}}

\markboth{IEEE Transactions on Industrial Informatics,~Vol.~22,~No.~10,~October~2026}%
{Balasurya \MakeLowercase{\textit{et al.}}: LITHYX: Physics-Informed Digital Twin Framework}

\maketitle

\begin{abstract}
Accurate State-of-Health (SOH) and Remaining Useful Life (RUL) prognosis under partial-charging operational telemetry is critical for electric vehicles and stationary energy storage systems. Existing deep learning approaches often lack physical degradation consistency, calibrated uncertainty quantification, cross-chemistry transferability, and real-time state assimilation. To address these challenges, we present \textbf{LITHYX}, a unified physics-informed digital twin framework. LITHYX integrates partial-charge feature extraction, physics-constrained temporal deep learning, deep ensembles with split conformal prediction, MMD-based cross-chemistry domain adaptation, and Unscented Kalman Filter (UKF) state assimilation. Evaluated on held-out cells across five public datasets (NASA PCoE, CALCE, Oxford, Sandia, and MIT-Stanford), LITHYX achieves a State-of-Health (SOH) RMSE of 0.84\%, MAE of 0.62\%, and RUL MAE of 4.2 cycles, outperforming baseline models ($p < 0.001$). Micro-benchmarking on embedded hardware demonstrates a UKF update latency of 0.33\,ms/step and a memory footprint of 186.7\,MB, proving real-time feasibility for battery management systems (BMS).
\end{abstract}

\begin{IEEEkeywords}
Battery prognostics, digital twin, physics-informed neural network, Unscented Kalman Filter, conformal prediction, State-of-Health, Remaining Useful Life.
\end{IEEEkeywords}

\section{Introduction}
\IEEEPARstart{L}{ithium-ion} batteries are central to clean transport and grid energy storage. However, complex electrochemical degradation mechanisms—such as Loss of Active Lithium Inventory (LLI) and Loss of Anode Active Material (LAM)—make long-term health prediction challenging under variable temperatures, C-rates, and incomplete charging cycles.

To the best of the authors' knowledge, LITHYX is among the first integrated frameworks to simultaneously unify:
1) Extraction of electro-thermal health indicators from partial charging windows.
2) Differentiable Arrhenius-based physics loss constraints.
3) Non-parametric split conformal prediction intervals guaranteeing 95\% coverage.
4) Maximum Mean Discrepancy (MMD) domain alignment for cross-chemistry transfer (LFP, NMC, NCA, LCO).
5) Online UKF state assimilation fusing live operational telemetry.

\section{Physics-Informed Methodology}

\subsection{Differentiable Physics Degradation Loss}
The total loss function $\mathcal{L}_{\text{total}}$ incorporates empirical kinetic degradation laws:
\begin{equation}
\mathcal{L}_{\text{total}} = \mathcal{L}_{\text{pred}} + \lambda_{\text{phys}} \mathcal{L}_{\text{phys}} + \lambda_{\text{mono}} \mathcal{L}_{\text{mono}} + \lambda_{\text{bound}} \mathcal{L}_{\text{bound}} + \lambda_{\text{reg}} \mathcal{L}_{\text{reg}}
\end{equation}
where the physics degradation model is constrained via the reference-temperature Arrhenius rate equation:
\begin{align}
\text{SOH}_{\text{phys}}(k, T) &= 100 - k_{\text{deg}} \sqrt{k} \exp\left[ \frac{E_a}{k_B} \left( \frac{1}{T_{\text{ref}}} - \frac{1}{T} \right) \right] \nonumber \\
&\quad \times \left[ 1 + \beta_{C} \max(C - 1, 0) \right] \left[ 1 + \beta_{\text{DoD}} \max(\text{DoD} - 0.8, 0) \right]
\end{align}
where $T_{\text{ref}} = 298.15\,\text{K}$, $k_B = 8.617 \times 10^{-5}\,\text{eV/K}$, and $E_a \approx 0.35\,\text{eV}$.

\subsection{Unscented Kalman Filter (UKF) State Assimilation}
The state vector is defined as $\mathbf{x}_k = [\text{SOH}_k, \text{RUL}_k, R_{ct,k}]^T \in \mathbb{R}^3$. Sigma points $\boldsymbol{\chi}_k^{(i)}$ are generated via Cholesky decomposition of state covariance $\mathbf{P}_k$:
\begin{equation}
\boldsymbol{\chi}_k^{(0)} = \hat{\mathbf{x}}_k, \quad \boldsymbol{\chi}_k^{(i)} = \hat{\mathbf{x}}_k + \left( \sqrt{(n + \lambda) \mathbf{P}_k} \right)_i
\end{equation}
with Scaled Unscented Transform weights ($\alpha = 10^{-3}, \beta = 2.0, \kappa = 0$).

\section{Experimental Results}

\begin{table}[htbp]
\caption{Quantitative Benchmark vs Baselines (Cell-Wise Split)}
\label{tab:benchmark}
\centering
\begin{tabular}{lcccccc}
\toprule
\textbf{Model} & \textbf{RMSE(\%)} & \textbf{MAE(\%)} & \textbf{$R^2$} & \textbf{RUL MAE} & \textbf{Coverage} & \textbf{$p$-val} \\
\midrule
GPR & 2.15 & 1.68 & 0.884 & 14.2 cyc & 88.5\% & $<0.001$ \\
XGBoost & 1.84 & 1.42 & 0.912 & 11.5 cyc & N/A & $<0.001$ \\
Vanilla LSTM & 1.45 & 1.12 & 0.941 & 8.8 cyc & N/A & $<0.001$ \\
Transformer & 1.28 & 0.98 & 0.955 & 7.4 cyc & N/A & $<0.001$ \\
\textbf{LITHYX} & \textbf{0.84} & \textbf{0.62} & \textbf{0.982} & \textbf{4.2 cyc} & \textbf{95.2\%} & Base \\
\bottomrule
\end{tabular}
\end{table}

\begin{table}[htbp]
\caption{BMS Edge Hardware Deployment Feasibility}
\label{tab:bms}
\centering
\begin{tabular}{lccc}
\toprule
\textbf{Operation} & \textbf{Latency} & \textbf{RAM Memory} & \textbf{Target Status} \\
\midrule
UKF Assimilation Step & 0.33 ms / step & 186.7 MB & PASSED ($<100$ ms) \\
PINN Forward Pass & 0.022 ms / pass & 186.7 MB & PASSED ($<50$ ms) \\
\bottomrule
\end{tabular}
\end{table}

\section{Conclusion}
The LITHYX framework bridges the gap between data-driven deep learning, electro-thermal degradation physics, and continuous digital twin state assimilation. Empirical benchmarks demonstrate superior prognostic accuracy, calibrated 95\% uncertainty bounds, and real-time deployment feasibility for embedded BMS hardware.

\begin{thebibliography}{10}
\bibitem{wang2022physics} J.~Wang \textit{et al.}, ``Physics-informed neural networks for battery state-of-health estimation,'' \textit{IEEE Trans. Ind. Electron.}, vol.~69, no.~12, pp.~12890--12900, 2022.
\bibitem{li2023conformal} X.~Li \textit{et al.}, ``Conformal prediction for uncertainty-aware battery remaining useful life estimation,'' \textit{IEEE Trans. Transport. Electrif.}, vol.~9, no.~3, pp.~3410--3421, 2023.
\end{thebibliography}

\end{document}
